\section{Effect-Safe Extraction}
\label{sec:effectful-extraction}

% Sources: ../egg-papers/2026-OOPSLA-eggcc/abstract.tex and 01-introduction.tex.
\sourcenote{Edited from \emph{Efficient Extraction for Effectful E-graphs},
  abstract and \S1; new connection to the thesis.}
Effect-safe extraction selects computations together with a valid order for their effects.
Prior approaches use integer linear programming (ILP),
  encoding program choices and effect-ordering requirements as constraints for a solver.
Solving this extraction problem can dominate optimization time.
Our work, published at OOPSLA 2026, introduces \emph{statewalk DP},
  a dynamic programming algorithm whose running time depends on a structural parameter,
  \emph{statewalk width}~\cite{eggcc}.
Its implementation in \texttt{eggcc} makes effectful extraction practical
  and enables useful domain-specific optimizations.

\subsection{Background: Effects and Control Flow}
\label{sec:background-effects}

% Source: ../egg-papers/2026-OOPSLA-eggcc/02-background.tex,
% Dataflow-based IRs and RVSDGs; project-specific discussion omitted.
\sourcenote{Edited from \emph{Efficient Extraction for Effectful E-graphs}, \S2.}
Compilers manipulate programs using \emph{intermediate representations} (IRs).
In a dataflow IR, nodes represent computed values and edges encode their dependencies.
Operations with side effects also have ordering dependencies:
  reordering memory writes or output can change behavior.
A \emph{state value} encodes this order, passing from one effectful operation to the next.
Each effectful operation consumes a state value and produces a new one,
  making its execution depend on the preceding operation.
An \emph{effect-safe} graph uses state values without duplicating or dropping them,
  so its effectful operations form a single sequence called a \emph{statewalk}~\cite{eggcc}.

% Sources: same background; 04-formal-rvsdgs.tex, purity and stateful values.
\sourcenote{Edited from \emph{Efficient Extraction for Effectful E-graphs}, \S\S2, 4.}
A term is \emph{effectful} if it produces a state, possibly alongside other values.
Values other than state values are called \emph{pure values},
  and a term producing only these values is called \emph{pure}.
Purity is a property of the value a term denotes, not of how that value is computed:
  a pure term can depend on a memory operation that produces its value.
Regionalized Value State Dependence Graphs (RVSDGs) combine these data and state dependencies
  with nested regions for control-flow constructs such as loops and conditionals~\cite{rvsdg}.
Each region has its own inputs, outputs, and statewalk.

\subsection{The Effect-Safe Extraction Problem}
\label{sec:effect-safe-problem}

% Sources: ../egg-papers/2026-OOPSLA-eggcc/03-overview.tex;
% 04-formal-rvsdgs.tex, semantics and authoring rules.
\sourcenote{Edited from \emph{Efficient Extraction for Effectful E-graphs}, \S\S3--4.}
An e-graph can contain equal terms that differ in whether their effects form a valid sequence.
Our semantics treats states as explicit values, allowing terms to duplicate or discard them;
  only the effect-safe terms can be used for code generation.
Rewrites preserve equality under this semantics,
  while extraction must also enforce the ordering constraint.
Congruence alone does not preserve effect safety:
  substituting an equal subterm can create two operations that consume the same state.

% Source: ../egg-papers/2026-OOPSLA-eggcc/03-overview.tex, Update example.
\sourcenote{Edited from \emph{Efficient Extraction for Effectful E-graphs}, \S3.}
Figure~\ref{fig:effect-safe-example} illustrates the problem with a memory operation,
  $\mathtt{Update}(a,x,\sigma)$.
It writes $x$ to address $a$ in state $\sigma$ and returns
  both the value previously stored there and the updated state.
The original program performs four updates, named $C_1$ through $C_4$:
\[
  \begin{aligned}
    C_1:\quad (v_1,\sigma_1)&=\mathtt{Update}(a,x,\sigma_0),\\
    C_2:\quad (v_2,\sigma_2)&=\mathtt{Update}(a,x,\sigma_1),\\
    C_3:\quad (v_3,\sigma_3)&=\mathtt{Update}(b,y,\sigma_2),\\
    C_4:\quad (v_4,\sigma_4)&=\mathtt{Update}(c,v_2,\sigma_3).
  \end{aligned}
\]
Here $a,b,c$ are memory addresses, $x,y$ are values, and $\sigma_0$ is the initial state.
The second update repeats the first, so $\sigma_2=\sigma_1$ and $v_2=x$.
Recording both equalities permits eliminating $C_2$ and passing $x$ directly to $C_4$.

% Source: ../egg-papers/2026-OOPSLA-eggcc/figures/overviewcombinedcropped.pdf,
% Figure 3; caption adapted from 03-overview.tex.
\begin{figure}[tbp]
  \sourcenote{Verbatim figure from \emph{Efficient Extraction for Effectful E-graphs},
    Figure~3; edited caption.}
  \centering
  \includegraphics[width=\linewidth,pagebox=cropbox]{effect-safe-extraction.pdf}
  \caption{A redundant update creates both effect-safe and effect-unsafe choices.
    (I) The original program.
    (II) The e-graph after adding the redundant-update equalities.
    (III) An effect-unsafe term: $C_2$ and $C_3$ consume the same state.
    (IV) An effect-safe term that removes $C_2$.
    Operators point to their operands; bold arrows track state dependencies.
    Bold outlines mark effectful terms.
    Index $[0]$ selects an update's old value and $[1]$ its new state.
    Gray shading and dotted lines mark equalities; red and green highlight
    the unsafe and safe choices, respectively.}
  \label{fig:effect-safe-example}
\end{figure}

% Source: same overview, effect-unsafe and effect-safe extractions;
% aliasing example made explicit.
\sourcenote{Edited from \emph{Efficient Extraction for Effectful E-graphs}, \S3;
  concrete aliasing case added.}
The e-graph also permits replacing $\sigma_2$ with $\sigma_1$
  while retaining $C_2$ to compute $v_2$.
Now $C_2$ and $C_3$ consume the same state, as in Panel~III.
Simply scheduling these operations in either order can change the result:
  if $a=b$ and $x\ne y$, executing $C_3$ before $C_2$
  makes $C_2$ return $y$ instead of $x$.
Panel~IV avoids this problem by selecting both equalities together,
  producing the statewalk $C_1,C_3,C_4$ and using $x$ for the last update.
Extraction must therefore coordinate its choices across e-classes.

% Source: ../egg-papers/2026-OOPSLA-eggcc/07-np.tex, hardness and fallback.
\sourcenote{Edited from \emph{Efficient Extraction for Effectful E-graphs}, \S5.}
We prove that deciding whether a given e-class contains any effect-safe extraction
  is NP-complete, even without a cost objective~\cite{eggcc}.
An optimizer can retain its input program as an effect-safe fallback,
  but its goal is to find an improved program among the alternatives.
Choosing a cheap term independently for each e-class can combine incompatible
  effect orders, and rejecting that combination can miss a valid extraction.

\subsection{Statewalk DP and Statewalk Width}
\label{sec:extraction-algorithm}

% Sources: ../egg-papers/2026-OOPSLA-eggcc/03-overview.tex;
% 06-tiger.tex, extractable sets and statewalk equivalence.
\sourcenote{Edited from \emph{Efficient Extraction for Effectful E-graphs}, \S\S3, 6.}
Our key insight is that different statewalks can enable the same pure computations.
For a statewalk, its \emph{extractable set} contains the pure e-classes
  from which a term can be built using only the effects on that statewalk.
Statewalk DP groups statewalks when their final terms belong to the same e-class
  and their extractable sets are equal.
Statewalks in the same group can replace one another when constructing further effect-safe terms.
In Figure~\ref{fig:effect-safe-example}, $\sigma_2=\sigma_1$ lets a statewalk repeat $C_2$
  any number of times.
These statewalks all enable the same pure e-classes.
Grouping them avoids enumerating all those repetitions.

% Source: same overview and 06-tiger.tex, algorithm description.
\sourcenote{Edited from \emph{Efficient Extraction for Effectful E-graphs}, \S\S3, 6.}
The algorithm uses a dynamic programming table to retain and reuse partial extractions.
Each entry is indexed by an effectful e-class and an extractable set,
  and stores a representative statewalk for that group.
Starting from the input state, the algorithm extends a statewalk with an effectful operation
  only when each of its pure operands can be extracted using that statewalk.
It computes the new extractable set and records the resulting group.
The implementation keeps a low-cost representative for each group,
  using a priority queue to explore cheaper choices first.
This preserves the compatibility information needed for future extensions
  while avoiding separate searches for interchangeable statewalks.

% Source: ../egg-papers/2026-OOPSLA-eggcc/06-tiger.tex, Complexity;
% 03-overview.tex, width-one example.
\sourcenote{Edited from \emph{Efficient Extraction for Effectful E-graphs}, \S\S3, 6.}
The \emph{statewalk width} $k$ is the largest number of distinct extractable sets
  for statewalks ending in any one effectful e-class.
The running example has width one, despite its unbounded number of statewalks.
For an e-graph containing $n$ explicitly stored terms, finding any effect-safe extraction
  takes $O(kn^2)$ time; prioritizing lower-cost choices adds a logarithmic factor.
Thus the feasibility problem is \emph{fixed-parameter tractable} in statewalk width:
  its running time is a function of $k$ times a polynomial in $n$
  whose exponent does not depend on $k$.
Width can be exponential in the worst case, but small width bounds how many
  different compatibility requirements the algorithm must track.

% Sources: ../egg-papers/2026-OOPSLA-eggcc/06-tiger-proofs.tex;
% 08-impl.tex, Cost Model. New scope note for report length.
\sourcenote{Edited from \emph{Efficient Extraction for Effectful E-graphs}, \S\S6--7;
  new scope note.}
For the core dataflow language, we prove that statewalk DP terminates,
  returns only effect-safe terms in the requested e-class,
  and finds one whenever one exists.
Its cost optimization is greedy and does not guarantee a globally minimum-cost program.
For space, this report gives the algorithm's main ideas;
  the full algorithm and correctness arguments appear in the paper~\cite{eggcc}.

\subsection{Implementation and Evaluation}
\label{sec:extraction-evaluation}

% Source: ../egg-papers/2026-OOPSLA-eggcc/08-impl.tex,
% Control Flow and Regions; The eggcc Program Optimizer.
\sourcenote{Edited from \emph{Efficient Extraction for Effectful E-graphs}, \S7.}
We implement statewalk DP in \texttt{eggcc}, which optimizes programs written in Bril,
  a simple imperative intermediate language.
It translates Bril to RVSDGs, applies \texttt{egglog} rules, extracts an effect-safe RVSDG,
  and translates the result back to Bril.
To handle loops and conditionals, extraction decomposes the e-graph into
  \emph{regionalized e-graphs}, each describing a single region with placeholders for nested regions.
It runs statewalk DP on these smaller problems and recursively fills the placeholders.

% Source: ../egg-papers/2026-OOPSLA-eggcc/10-eval.tex, Setup and RQ1;
% values from nightly/jscomputedmacros.tex, checked against main-camera-ready.pdf.
\sourcenote{Edited from \emph{Efficient Extraction for Effectful E-graphs}, \S8.}
We compare statewalk DP with ILP extraction under the same cost model.
Across 96 benchmarks producing 9,610 regionalized e-graphs,
  statewalk DP completes 92\% of extractions in under a millisecond and all within 61.9\,ms.
Its geometric-mean speedup over the ILP solver Gurobi is $520\times$,
  counting timeouts as five minutes and using reported solving times for infeasible instances.
These figures measure extraction search, excluding the overhead of transferring and preparing e-graphs.

% Sources: same evaluation, statewalk-width measurements;
% 08-impl.tex, liveness and AC mitigation.
\sourcenote{Edited from \emph{Efficient Extraction for Effectful E-graphs}, \S\S7--8.}
Statewalk width is one in 94.7\% of the measured Bril and PolyBench regions
  and at most five in 98.4\%.
Two outliers reach width 106,259 because independent memory reads admit many orderings.
Dropping pure values that no future operation needs and restricting redundant read orderings
  reduce the largest table for an e-class in the highest-width case to 24 entries.
The parameter thus explains expensive instances and guides practical improvements.

% Sources: same evaluation, RQ2; 10-eval-fenwick.tex, RQ3;
% numerical speedup from nightly/jscomputedmacros.tex.
\sourcenote{Edited from \emph{Efficient Extraction for Effectful E-graphs}, \S8.}
Where ILP extraction finishes, statewalk DP produces programs with similar measured runtimes.
The prototype remains slower than the LLVM compiler on the benchmark-suite aggregates,
  but it supports useful optimizations tailored to an application.
For a Fenwick tree, a data structure for maintaining cumulative sums under updates,
  we add three rules that replace a loop computing the lowest set bit of a positive integer $n$
  with the bitwise expression $n\mathbin{\&}(-n)$.
In this case study, the resulting program runs $6.52\times$ faster than LLVM with full
  \texttt{-O3} optimization.
These rules combine with memory and control-flow transformations to improve the executable program.

\sourcenote{New connection to the proposed verified compiler.}
Effect-safe extraction constructs candidates for the proposed compiler
  in Section~\ref{sec:verified-compilation}, which must validate their effect ordering.
It assumes the e-graph's equalities come from sound rewrites;
  its correctness argument does not verify the rewrite rules or the complete compiler.
Combining that guarantee with proofs of rewrites and analysis results,
  and a verified checker, remains part of the proposed work.

\subsection{Related Work}
\label{sec:related-effects}

% Sources: ../egg-papers/2020-big-paper/09-related.tex, Superoptimization;
% ../egg-papers/2026-OOPSLA-eggcc/11-related.tex, Effect-safe Extraction in Peggy.
\sourcenote{Edited from the related work of \emph{egg}
  and \emph{Efficient Extraction for Effectful E-graphs}.}
Peggy uses e-graphs to optimize programs with loops and conditionals,
  then selects a program through constraint solving~\cite{peggy}.
Translating that result back to a control-flow graph, which records possible transfers
  of execution between instruction blocks, is a separate obligation~\cite{peg-cfg}.
Our work makes effect safety an explicit extraction constraint and studies its complexity.

% Source: ../egg-papers/2026-OOPSLA-eggcc/11-related.tex, Pure Extraction in Practice.
% Julia comparison checked against arXiv:2502.17075v2, Section 3.3.
\sourcenote{Edited from \emph{Efficient Extraction for Effectful E-graphs}, related work;
  Julia comparison updated from its paper.}
Cranelift's e-graph design preserves effectful operations in the surrounding control-flow graph~\cite{aegraph}.
The Julia optimizer likewise preserves that structure, while allowing marked rewrites to remove effectful calls~\cite{eqsat-julia}.
These approaches restrict which effects extraction must reconstruct;
  ours searches the e-graph for a valid effect ordering.

% Source: ../egg-papers/2026-OOPSLA-eggcc/11-related.tex, Pure Extraction in Practice;
% 01-introduction.tex, parameterized complexity discussion.
\sourcenote{Edited from \emph{Efficient Extraction for Effectful E-graphs},
  introduction and related work.}
Even without effects, extraction becomes difficult when the cost model accounts for sharing.
A cheap expression tree need not yield a cheap directed acyclic graph
  when repeated computations are counted only once.
Goharshady and colleagues give an optimal extraction algorithm
  parameterized by e-graph \emph{treewidth}, a measure of how closely
  the graph's connectivity resembles that of a tree~\cite{goharshady}.
Sun and colleagues connect extraction to Boolean circuits
  and obtain another algorithm parameterized by treewidth~\cite{glenn-treewidth}.
These results show how structural parameters can make an otherwise hard extraction problem tractable.
Our work applies this strategy to a different obstacle:
  finding any effect-safe extraction, parameterized by statewalk width~\cite{eggcc}.
